\begin{note}
    Для доказательста теоремы нужна теория о сопряжённых операторах~(\ref{conjugate-op-section}). На лекции это было до теоремы.
\end{note}

\begin{lemma}\label{8.4_lemma}
        $E$~--- банахово пространство,
        $A \in \L(E)$, $A$~--- ограниченный снизу. 
        Тогда $\overline{\im A} = \im A$.
\end{lemma}
\begin{proof}
    Пусть \(\left\{y_n \right\} \subset \im A\), такая что \(y_n \rightarrow y\) по норме. Представим как $y_n = A z_n, z_n \in E$. Последовательность $\{y_n\}$ сходится, а значит она фундаментальна.
    \[ \epsilon > \|y_{n+p}-y_n\|=\|Az_{n+p}-Az_n\| \geq m \cdot \|z_{n+p} - z_n \|.\] 
    Итак, $\{z_n\}$ - фундаментальна в $E$~--- банаховом, откуда $\exists \lim z_n = z$. Следовательно, $y_n = A z_n \to A z = y$. То есть $y \in \im A$, а значит $\im A$ замкнут.
\end{proof}

\begin{theorem}[Банаха об обратном операторе (8.4) в случае гильбертовых пространств]\label{8.4}
     Пусть $H(\K)$~--- гильбертово простраство, $A \in \mathcal{L}(H)$,
     $A$~--- биекция $H$ на $H$. Тогда $A^{-1} \in \mathcal{L}(H)$.
\end{theorem}

\begin{proof}  % доказательство, которое невозможно забыть
    \begin{enumerate}
        \item Докажем, что $(A^*)^{-1} \in \mathcal{L}(H)$. По теореме~\nameref{8.1} для этого достаточно показать, что
        $A^*$~--- ограниченный снизу, то есть 
        $\exists m > 0:\ \forall x \ \|A^* x\| \geqslant m \|x\|$.
        \item Если неравенство $\|A^* x\| \geqslant m \|x\|$ верно для $x$, то оно также верно для $\alpha x$ ($\alpha \neq 0$). С учётом этого можно сделать трюк.

        \textbf{Трюк:} рассмотрим множество 
        $S = \{x \in H \: | \: \|A^*x\| = 1\}$. Следовательно, для выполнения $\exists m > 0:\ \forall x \ \|A^* x\| \geqslant m \|x\|$ достаточно показать, что
        \[\exists m > 0\!:\: \left(\forall x \in S \; 1 \geqslant m \|x\| \right) \Longleftrightarrow \exists m > 0\!:\: \left(\forall x \in S\: \|x\| \leqslant \frac{1}{m} \right)\]
        То есть надо доказать, что $S$~--- ограниченное множество.
        % Греция -- штука странная. Куда плавал Одиссей?
        % Будете себя плохо вести -- докажем её ещё раз
        Покажем, что $S$~--- слабо ограниченное, то есть
        \[
        \forall y \in H\; \exists K_y\; \forall x \in S \left(|(x, y)| \leqslant K_y \right)
        \]
        
        тогда по теореме~\nameref{han} будет следовать ограниченность $S$. \(y\) играет роль функционала, так как $H$ гильбертово, по теореме \nameref{Riss-Freshe}, его действие можно заменить на скалярное произведение с элементом $H$.

        % главное место в доказательстве
        Используем, что $A$~--- биекция $H$ на $H$: $\forall y \in H \; \exists z \in H \; Az = y$

        \[ |(x, y)| = |(x, Az)| = |(A^*x, z)| \leq \underbrace{\|A^* x\|}_{=1\text{ на $S$}} \cdot \underbrace{\|z\|}_{=K_y\text{ (обозначим)}}\]
        Действительно слабо ограничено, а значит и просто ограничено и оператор $A^*$ ограничен снизу.
        \item Итого по теореме~\nameref{8.1}, $\exists (A^*)^{-1} \in \L(\im A^*)$. 
        Из биективности $A$, $\ker A = \{0\}$, $\im A = H$, по формуле теоремы~\nameref{9.2} имеем 
            $\ker A^* = \{0\}$ и $\overline{\im A^*} = H$.
        По лемме~\ref{8.4_lemma} $\im A^* = \overline{\im A^*} = H$. Имеем инъективность и сюрьективность, то есть $A^*$~--- биекция.
        
        Выходит что $A^*$ удовлетворяет условиям теоремы и $\exists (A^*)^{-1} \in \L(H)$.
        Провернём рассуждения ещё раз для $A^*$, получим $(\underbrace{A^{**}}_{=A})^{-1} \in \L(H)$~--- перешли к исходному оператору $A$. 
    \end{enumerate}
\end{proof}